\documentclass{article}

% Language setting
% Replace `english' with e.g. `spanish' to change the document language
\usepackage[english]{babel}

% Set page size and margins
% Replace `letterpaper' with `a4paper' for UK/EU standard size
\usepackage[letterpaper,top=2cm,bottom=2cm,left=3cm,right=3cm,marginparwidth=1.75cm]{geometry}

% Useful packages
\usepackage{amsmath}
\usepackage{amssymb}
\usepackage{graphicx}
\usepackage[colorlinks=true, allcolors=blue]{hyperref}

\title{Lattice QCD Notes}
\author{Shivam Panda}

\date{}
\begin{document}
\maketitle

\section{Background}
Typically, the standard Hamiltonian is the generator of evolution. So, in typical QM, the transition amplitudes (the probability of a particle starting at $x_i$ at $t_i$ to be found at $x_f$ at $t_f$ is given by: 
\begin{equation}
    <x_f|e^{-iH(t_f-t_i)}|x_i>
\end{equation}
However, the imaginary exponential leads to computational troubles due to rapid oscillations, so a Wick Rotation is performed with transforms $it \xrightarrow{} t$. This Wick Rotation switches Minkowski space to Euclidean space, where time behaves as another spacial dimension and allows the computation using a 4D grid. However, this change in time causes significant limitations on simulating dynamics and time-evolution in the lattice. This leads to exponential decrease rather than oscillations, which is much more computationally stable. With this transformation, the transition amplitudes can be written as: 
\begin{equation}
    <x_f|e^{-H(t_f-t_i)}|x_i> =\int Dx(t)e^{-S[x]}
\end{equation}
This a 1-dimensional path integral, where $Dx(t)=A\int dx_1dx_2dx_3...dx_{N-1}$ as $N \xrightarrow{} \infty$. Each $dx_i$ represents all the possible positions at different slices of time. Essentially, it is an integral of all possible functions of $x(t)$. Paths with higher action are exponentially supressed, which means that the leading contributor is the path with minimal action.
After the Wick Rotation, the action and Lagrangian change
\begin{equation}
    S[x] = \int dtL(x,\dot{x})=\int dt[\frac{m\dot{x}^2}{2}+V(x(t))]
\end{equation}
The most significant difference is the positive sign in the Lagrangian. Using this formulation, the propagator can be written by resolving the identity into energy eigenstates. Since the Hamiltonian is diagonal with the energy eigenstates, the exponential simplifies from the Hamiltonian to simple energy eigenvalues.
\begin{equation}
    <x|e^{-HT}|x>=\sum <x|E_i>e^{-E_iT}<E_i|x>
\end{equation}
Due to the exponential suppression, as $T \xrightarrow{} \infty$, only the ground state remains and $<x|e^{-HT}|x>=e^{-E_0T}|<x|E_0>|^2$. Since $\psi_n(x)=<x|E_n>$, which is just an energy eigenstate written the position basis. The normalization properties of the wavefunction lead to $\int |\psi_n(x)|^2dx=1$, so that the ground state can be extracted by integrating all x.
\begin{equation}
    \int dx <x|e^{-HT}|x> =e^{-E_0T}
\end{equation}
However, this equation only holds true as $T \xrightarrow{} \infty$
\subsection{Discrete Path Integral}
Since the computer cannot simulate all possible functions x(t), it is possible to make the t-axis discrete, and describe the position $x(t_i)$, and integrate with all possible x to replicate the idea of integrating over all possible functions. Therefore, there are as many integrals as there are grid points on the t-axis. 
The t-axis can be made discrete by $t_j=t_0 + ja$, where $a = \frac{t_f-t_0}{N}$. With this, a path described as a vector $x = \{x(t_0),x(t_1), \dots ,x(t_N)\}$. Since the end points are fixed, the path integral becomes $Dx(t)=A\int dx_1dx_2dx_3...dx_{N-1}$ as $N \xrightarrow{} \infty$. A is simply a normalization factor. Using this, the action integral can be determined using discrete approximations of the derivative.
\begin{equation}
    \int_{t_j}^{t_{j+1}}dt =\int_{t_j}^{t_{j+1}}a[\frac{m}{2}(\frac{x_{j+1} - x_{j}}{a})^2+\frac{V(x_{j+1})+V(x_j)}{2}]
\end{equation}
The derivative is replaced by the limiting definition of the derivative, where dt is taken as a, and the potential is determined at an interval by taking the average value. Using all of this, the propagator and path integral can be computed.
\begin{equation}
    <x|e^{-HT}|x>=A\int_{-\infty}^{\infty}dx_1dx_2\dots dx_{N-1}e^{-S_{lat}[x]}
\end{equation}
\begin{equation}
    S_{lat}[x]=\sum_{j=0}^{N-1}[\frac{m}{2a}({x_{j+1}-x_{j})^2 + aV(x_j)]}
\end{equation}

\subsection{Monte Carlo Evaluation}
In the previous section, the probability amplitude was a path integral, however the every path had the same boundary conditions.
However, in this next section, the path integrals will be evaluated at every possible boundary condition, so it is effectively an integral of path integrals in some sense.
To find the correlation between $x(t_2)$ and $x(t_1)$, 
\begin{equation}
<<x(t_2)x(t_1)>> = \frac{\int Dx(t) x(t_2)x(t_1)e^{-S[x]}}{\int Dx(t) e^{-S[x]}}
\end{equation}
Here, both path integrals are integrating over all paths for a set of boundary conditions, and then also integrating over all possible boundary conditions.
In general, the $<<x(t_1)>>$ gives information about on average, what is the position of the particle at $t_1$, where the average is taken over all paths.
The correlation, like shown in Eq (9), gives information about the correlation between the position of the particle at $t_1$ and $t_2$. If $t_1 = t_2$, then 
this would give ifnromation about the variance of $x(t_1)$, which is essentially the width of the position due to quantum fluctuations.
Since the integral is over all possible boundary conditions, the denominator can be written as
\begin{equation}
    \int Dx(t) e^{-S[x]} = \int dx <x|e^{-HT}|x> = \sum e^{-E_n*T}
\end{equation}
This integral over boundary condition is the trace of an operator, which is just the sum of its eigenvalues. 
Going from this integral to sum is simply changing the basis for this trace calculation, moving from position to energy basis.
This equation is considered the partition function for a system.
The numerator can be simplified by considering that $t_f - t_1 = T$ and $t = t_2 - t_1$.
A similar approach can be taken with the numerator of Eq 9, but more care needs to be taken. \newline
\begin{equation}
    \sum_n |E_n>e^{-H(t_f - t_2)}\hat{x}e^{-H(t_2 - t_1)}\hat{x}e^{-H(t_1 - t)}<E_n|
\end{equation}
This expression rewrites the operator in energy basis, but it also gives intuition as to what the operator actually does.
Essentially, it allows the system to evolve from $t_i$ to $t_1$, then applies the position operator. 
Then, allows the system to evolve from $t_1$ to $t_2$, and then applies another position operator at $t_2$. 
Finally, it allows the system to evolve for the rest of the time considered.
By switching to the energy basis, the Hamiltonian becomes diagonal, so the evolution terms simplify.
\begin{equation}
    \sum_n e^{-E_n(t_f - t_2 + t_1 - t_i)} <E_n|\hat{x}e^{-H(t_2 - t_1)}\hat{x}|E_n>
\end{equation}
\begin{equation}
    \sum_n e^{-E_nT} <E_n|\hat{x}e^{-(H - E_n)t}\hat{x}|E_n>
\end{equation}
Using this equation, and assuming that $T >> t$, where T is very large, the correlation
\begin{equation}
    G(t) = \frac{e^{-E_0T<E_0|\hat{x}e^{-(H - E_0)t}\hat{x}|E_0>}}{e^{-E_0T}} = <E_0|\hat{x}e^{-(H - E_0)t}\hat{x}|E_0>
\end{equation}
Since, as T becomes very large, the ground state dominates. Using parity arguments, the position operator is odd, while the $|E_0>$ is even.
Therefore, as t is large, 
\begin{equation}
    G(t) = |<E_0|\hat{x}|E_1>|^2 e^{-(E_1 - E_0)}t
\end{equation}
Using this formulation, the first energy eigenstate can be determined by evaluating the correlation function at t and t+a.
\begin{equation}
    log(\frac{G(t)}{G(t+a)}) = (E_1 - E_0)a
\end{equation}
$<E_0|\hat{x}|E_1>$ the transition matrix element that physically determines the probability of the position operator kicking the first energy eigenstate into the ground state.
Using this technique, $<<\Gamma[x]>>$ can be determined by generating a large number of configurations (paths) on the grid such that the probability
of obtaining a particular path is proportional to $e^{-S[x]}$, and thus, the expected value of $\Gamma[x]$ is simply a weighted average of these configurations.
This is a Monte Carlo estimator for the true expected value of the functional, however the uncertainty decreases with $1/\sqrt{N}$. \newline
Generating paths with this probability proportional to the exponent of the action is not a simply task however it can be performed by the Metropolis Algorithm.
\begin{equation}
    <<\Gamma[x]>> = \frac{1}{N_{cf}} \sum_{\alpha}^{N_{cf}} \Gamma[x^{(\alpha)}]
\end{equation}
This begins by start with a random path $x^0$. Then to generate $x^1$, all of the sites of the lattice are visited and the values at those sites are randomized by this procedure.
\begin{list}{}{}
    \item 1) Generate a random number ($\alpha$) $\in [-\epsilon, \epsilon]$
    \item 2) Replace $x_j$ with $x_j$ + $\alpha$
    \item 3) If action is reduced, actually replace $x_j$ with this new value
    \item 4) If action  is increased, then generate a random number ($\beta$) between 0 and 1. Retain new value if $e^{-\Delta S} > \beta$. Otherwise, keep old value and move to next site.
\end{list}

The value of $\epsilon$ is critical since if it is too high, then most new values for $x_j$ will be rejected, while if it is too small, many will get accepted, but the values will not change significantly.
The goal is that $x_j$ should pass with probability 0.4 - 0.6.
Additionally, paths should be systematically chosen as to eliminate the correlation between successive paths.
Lastly, the first 5N or 10N should be disregarded since it is extremely random and useless.

\subsection{Statistical Errors}
It is incredibly computationally intensive to estimate the statistical uncertainty of these monte carlo estimators by repeating the entire estimate a 100 times. Therefore, there is Statistical Bootstrapping.
Essentially, from the given values of G(t) for each path, a random selection of these G(t)'s can be picked to find an estimate, and these estimates can be used to determine the statistical errors.
Another useful process is called binning, where a systematic sampling of configurations and their corresponding G(t) values can be taken to find estimates and therefore be partially averaged to get estimates. 
These estimates can then be used to determine the statistical errors.
One key limitation of binning is that within a bin, the configurations could possibly be correlated, and therefore, the errors estimated could be much lower than the true errors.
This would occur if $N_corr$ is taken to be too small.

\section{Field Theory on Lattice}
With transition from QM to Field Theory is essentially adding spatial dimension. The $x(t)$ from before could be considered a field theory of 0 spatial and 1 temporal dimension.
Instead of have a ground state energy eigenstate, the field has a ground state $|0>$, which corresopnds to no particles. The operators $\phi$ behaves like the position operator, 
and it is possible for this operator to excite the field above its ground state, and create a particle. \newline
The nodes of the 4D Euclidean grid, separated by the lattice spacing a, all have corresponding values of the field
The partition function still exists in path integral form: 
\begin{equation}
    Z = \int e^{-S[\phi]} \prod d\phi(x_j)
\end{equation}
This is essentially a path integral over all possible field configurations.
$S[\phi]$ corresponds to the field action. The spatial and temporal derivates transform in the differences of the scalar field with neighboring field values.
Similar to the operators used in QM to study excitations from the ground state, operators like: 
\begin{equation}
    \Gamma(t) = \frac{1}{\sqrt{N}} \sum \phi(\vec{x_j}, t)
\end{equation}
Summing over all the spatial vectors enforces zero 3-momentum. Since the excitations correspond to particle creation, their energies are teh masses of the particles.
Any non-zero momentum would deconstructively interfere and cancel itself out since it would repeat throughout the lattice. (QUESTION)
\begin{equation}
    <<\Gamma(t)\Gamma(0)>> = |<0|\Gamma(0)|\phi : \vec{p} = 0>|^2 e^{-m_\phi t}
\end{equation}
Where $|\phi : \vec{p} = 0>$ represents a single $\phi$-particle state with zero three momentum and $m_\phi$ represents the mass of the particle. \newline
Often, additional terms are added into the Lagrangian to add local interactions to account for the ultraviolet cutoff.
The ultraviolet cutoff is simply the lattice property that there is a minimum wavelength ($\lambda = 2a$), where a is the lattice spacing.
Sometimes, these additional terms influence new modes of oscillations and solutions which are considered "numerical ghosts".

To improve these discretization errors, it is possible to make the transformation
\begin{equation}
    x_j = \tilde{x_j} + \psi_1 a^2 \Delta^(2)\tilde{x_j} + \psi_2 a^2 \omega_0^2 \tilde{x_j}
\end{equation}
With the correct choices of $\psi_1$ and $\psi_2$, all $O(a^2)$ terms can cancel in the action, and the discretization errors will be $O(a^4)$ and all numerical ghosts can be removed.
\section{QCD on a Lattice}
\subsection{Classical Gluons}
For some background, the vector potential and the field tensor are very different in QCD.
In QED, the vector potential's components are simply scalars, but in QCD, they are 3x3 Hermitian Traceless matrices.
Each component of the vector potential can be viewed as a linear combination of the 8 Gell-Mann matrices, each of which correspond to a different gluon.
Since the vector potential is a linear combination of these matrices, each of the "coefficients" of this linear combination is a 4-vector itself, but with scalar components.
To understand what these scalar components are doing, it is important to first understand what the Gell-Mann matrices actually represent, and what QCD itself is.
Just like the vector potential in QED is the "correction" term for local U(1) changes to the wave function, the vector potential for QCD represents corrections to local changes of the color.
In QCD, color is represented as a $\mathbb{C}^3$, which is a 3-component complex vector. Each of these components correspond to the 'probability' of color red, blue, or green.
The fundamental concept of QCD is that the universe is invariant to local SU(3) changes, which are essentially rotations in this $\mathbb{C}^3$ space that switch around the probability amplitudes of each color while preserving the total probability.
Therefore, the vector potential exists to be the "correction" term for these lcoal SU(3) changes, which is why the Gell-Mann matrices are themselves SU(3) matrices, so each gluon corresponds to a different type of rotation in $\mathbb{C}^3$. 
These 8 gluons serve as the basis of any general rotation in this space, and therefore each component of the 4-vector potential corresponds to a linear combination of these basis SU(3) matrices. \newline

In QED, the vector potential essentially changes the phase of the wave function of the charged particles as the correction to the local U(1) changes, and in QCD, the vector potential changes the orientation of the color vector in $\mathbb{C}^3$.
However, there is an interesting subtlety to this. In QED, because U(1) is abelian, the vector potential effects simply sum up, so the phase change of two vector potentials overlapping is simply the sum of them. This is because 2D rotations commute.
In QCD, the SU(3) matrices do not commute, and so in general, overlapping vector potentials do not simply sum. This is what fundamentally leads to gluon-gluon interaction. 
When these Gell-Mann matrices (the basis gluons) act on a vector in $\mathbb{C}^3$, their contibutions do not simply sum up, and the fact that their commutators are non-zero is what leads to gluons carrying color and self-interacting.
This is the reason why the field tensor is defined as the 4D curl the vector potential and the commutator between the various components of the 4-vector potential. \newline

The components of the scalar 4-vector coefficients for each Gell-Mann matrix has an interesting physical relevance. Essentially, if the 4-vector coefficient has a non-zero time-like component, then the color vector at that given point will essentially precess and change over time.
If the space-like components of the scalar 4-vector coefficients is non-zero, as the quark moves through space with some velocity, the color vector will change.
This behaves very similar to an E-field and B-field, but instead of the phase accumulation of a charged particle, the color vector changes for a quark.
Keep in mind that each component of the field tensor actually is a 3x3 Hermitian matrix that corresponds to a SU(3) transformation. To understand the need for the commutator, imagine a moving around a tiny square with a non-zero 4-vector potential.
Suppose that the square is in the xy-plane.
 It first feels the effect of the $A_x$ component and then $A_y$, and then the reverse. 
 Because these effects do not commute, the field tensor needs to account for this.
Besides that, the field tensor is exactly measuring the twisting (or 4D curl) of the 4-vector potential. This is essentially like the parallel transport of general relativity, where when the parallel transport is non-zero (going around a full loop doesn't bring you back to where you start), that is a sign of curvature of the space. 
Just like that, here, the parallel transport (the full loop of the quark) will not return to where it started because of the non-abelian nature of SU(3).
Suppose that the quark first travels in x and then y. It would feel the 4-vector potential's x-comopnent transformation and then y, and because the order matters, the commutator must be there.
There is an intrinsic rotation of the 'axes' themselves due to the non-abelian nature of SU(3), which is the intrinsic rotation. Additionally, there is further rotation due to the 4D curl of the 4-vector potential itself. The net twisting is a combination of these, which is what is captured by the field tensor. \newline

Each component of the field tensor represents the net rotation in $\mathbb{C}^3$ of making an infinitesmal loop in some plane. For example, the $F_{\mu \nu}$ represents the total SU(3) rotation of making a loop in the $\mu - \nu$ plane.
The lagrangian contains the sum of all of the traces of the square of the components. Fundamentally, $F_{\mu \nu}$ represents the energy or work required to maintain that curvature in color space. Taking the trace and summing all of them correspond to the total work required to maintain some curvature. 
This is why the lagrangian is essentially a measure of the total cost of all of these twists in the 4-vector potential.
In the lagrangian, there are really 3 distinct terms. There can be massless gluon waves, which behave exactly like the photon in QED. Then there are 3-gluon interactions, which is essentially energy of 1 gluon splitting into 2 other gluons. Then, there is the 4-gluon vertices, which are 2 gluons scattering off of one another. These interaction terms are what give the strong force its sticky nature.

To maintain the SU(3) guage invariance on a discrete lattice, instead of tracking the vector potential, there is a link object, which is related to the 4-vector potential.
Additionally, the way that the 4-vector potential is defined, it is the instantaneous color rotation that the quark would experience. However, the actual rotation between nodes of the grid is this SU(3) matrix exponentiated.
After exponentiating, the U corresponds to a finite SU(3) transformation rather than the infinitesmal transformation that the 4-vector potential represents.
\begin{equation}
    U_{\mu}(x) = P(exp({-i \int_{x}^{x+a\mu}}{g A_{\mu} dy}))
\end{equation}
The Path Ordered integral is a necessity for dealing with non-abelian multiplication.
In practice however, because the 4-vector potential can be assumed to be approximately constant through this integration region: 
\begin{equation}
    U_{\mu}(x) \approxeq exp(-igaA_{\mu})
\end{equation}
Now, these $U_{\mu}$ are considered the fundamental tools on the grid, not the 4-vector potential, since they actually correspond to the vertices between nodes.
It is important to keep in mind that on the grid, the hermitian conjugate of some $U_{\mu}$ corresponds to the opposite direction vertex at the exact same node. 
Therefore, considering a 2D slice of space time, really the only vertices that are significant are $U_x$ and $U_y$, which correspond to the positive directions of x and y. The negative can be found by simply taking the hermitian conjugate.

A significant quantity that can be extracted from the $U_{\mu}$ is the trace of its products along a closed loop. First of all, by the structure of it, it is guage invariant. 
Additionally, the diagonal elements of each $U_{\mu}$ gives information of how much of the red-color stays the red-color or green with green. Therefore, taking the trace of any given $U_{\mu}$, it gives information of how much of the color stays the same between two nodes. 
Examining the trace of a closed loop gives information of the difference between the color before and after going through the entire loop, and so it gives a measure of the amount of curvature (twisting) of the 4-vector potential and thereby the $U_{\mu}$.
Therefore, the Wilson Loop function has the structure: 
\begin{equation}
    W(C) = \frac{1}{3}Tr(P(exp(-i \int gAdx)))
\end{equation}
Where the integral is along the closed loop. This simplifies into: 
\begin{equation}
    W(C) = \frac{1}{3}Tr(U_{\mu}(x)U_{\nu}(x + a\mu)\dots U^{\dag}_{\nu}(x))
\end{equation}
It is clear that since the loop ends with $U^{\dag}_{\nu}$ and starts with $U_{\mu}$, the loop starts with going to the right and ends with coming down at the same node, which is why it is a closed loop.
The reason for the factor of a third is because the null case is if there is no parallel transport, and so the resulting matrix multiplication is just I, which has a trace of 3, since this is a 3x3 complex matrix. \newline

A useful operator is the plaquette operator, which is simply the Wilson Loop operator around the smallest possible loop on the grid.
With a very weak classical $A_{\mu}$, the $P_{\mu \nu} \approx 1$, since there is very little curvature, and the twisting is negligible.
Therefore, the tension or energy cost of this non-zero vector potential is proportional to the deviation of the plaquette from 1. 
\begin{equation}
    S = \beta \sum (1 - P_{\mu \nu})
\end{equation}  
By taylor expanding the plaquette operator and using Stoke's theorem in 4D.
\begin{equation}
    S = \int d^4x \sum_{\mu \nu}{\frac{1}{2}Tr(F^2_{\mu \nu}) + \frac{a^2}{24}Tr(F_{\mu \nu} (D^2_{\mu} + D^2_{\nu})F_{\mu \nu} + \dots)}
\end{equation}
The higher order terms in a are essentially compensating for discretization errors. $D_{\mu}$ represents the covariant derivative, which essentially accounts for the rotation of the axes themselves due to the non-abelian SU(3).
It is possible to remove discretization errors by account for more points on the grid. For example, the rectangle operator $R_{\mu \nu}$ is essentially the plaquette operator but on a 2a x a rectangle rather than an a x a square.
Using this rectangle operator and the plaquette operator, it is possible to take a linear combination to only have errors of order $a^4$.
\begin{equation}
    S = -\beta \sum \frac{5P_{\mu \nu}}{3} - \frac{R_{\mu \nu} + R_{\nu \mu}}{12}
\end{equation}
The action above only has errors of order $a^4$.
\subsection{Quantum Gluons}
In the qauntum field theory analysis, there is the idea of tadpoles. Essentially, they are a loop in the Feynmann diagram and represent quantum fluctuations of a particle being created and annihalted instantaneously. 
This creates a significant amount of noise in these results, and so the values must be renormalized so that the noise is filtered out.
This renormalization is done by a factor $u_0$, which is essentially the mean link value. Since the plaquette has 4 links, it must be renormalized by $u^4_0$, and the rectangle operator must be renormalized by $u^6_0$.
Essentially, these tadpoles are gluons that are popping in and out of existence in the vacuum, and so for each link, $u_0$ represents the amount of artificial noise that is being created and must be renormalized.
One possible way to measure $u_0$: 
\begin{equation}
    u_0 = <0|P_{\mu \nu}|0>^(\frac{1}{4})
\end{equation}
This is essentially the expected fluctuations of the plaquette operator in a vacuum, and this corresponds to the quantum fluctuations of the gluons in a vacuum. \newline

In addition to these tadpole renormalizations, another renormalization must be made for unphysical contributions for effects that have higher wavelength than can exist on the grid. 

\subsection{Gluon Path Integrals}
The Metropolis algorith can be used, but now instead of randomly varying the field value across space and time, one must vary this $U_{\mu}$ parameter.
$U_{\mu}$ is a 3x3 complex matrix, so when it is varied, it is essentially being randomly multiplied by a random SU(3) matrix from a chosen set.
\begin{equation}
    U_{\mu} \rightarrow MU_{\mu}
\end{equation}
The set of SU(3) matrices M can be generated by first creating a set of hermitian matrices, and then scaling it down and adding $1 + i\epsilon H$. Again, $\epsilon$ determines the size of how much the $U_{\mu}$ actually updates throughout the Metropolis algorithm.

\end{document}